\documentclass[11pt]{letter}
\usepackage[margin=1in]{geometry}
\usepackage{parskip}
\usepackage[hidelinks]{hyperref}
\signature{Haitham A. El-Ghareeb}
\address{Information Systems Department\\Faculty of Computers and Information Sciences\\Mansoura
University, Mansoura, Egypt\\\href{mailto:helghareeb@mans.edu.eg}{helghareeb@mans.edu.eg}}
\begin{document}
\begin{letter}{The Editors\\\textit{Scientific Reports}}
\opening{Greetings,}

I am submitting the revised version of manuscript 2650209c-ff1e-4df4-aeb0-75ef3b35e564,
``Aggressive feature reduction degrades the selective reliability and clinical net benefit of
clinical risk models,'' for further consideration as an Article in \textit{Scientific Reports}. A
point-by-point response to both reviewers accompanies this letter.

\textbf{What the paper reports.} Feature reduction is a near-universal preprocessing step in
clinical risk modelling and is conventionally justified by showing that discrimination (AUROC)
changes little. This study asks whether that justification is sufficient. It introduces a
pre-registered, leakage-safe protocol certifying a reduced model on the properties a deployed risk
score is actually relied upon for --- probability calibration, selective (risk--coverage)
reliability, clinical net benefit, distribution-free conformal informativeness, and between-subgroup
equity --- as a function of the retained-feature budget. Across twelve public clinical datasets,
three learners and three feature-ranking methods (108 configurations, 30 repetitions each),
aggressive reduction degrades selective reliability, clinical net benefit and conformal
informativeness at least as much as discrimination, and in a subset of settings with no
discrimination warning at all, while probability calibration is mixed. At a matched budget of a
single retained feature every dataset is materially harmed, and on the mammographic-mass cohort
reduction flips the biopsy decision for one patient in five at a clinically conventional threshold.

\textbf{What changed in revision, including where the revision went against us.} Reviewer~1
identified that our screening criterion was confounded by the budget-rounding rule. Rather than
patch the affected analyses I re-ran the entire experimental grid under a new hash-pinned
configuration, added an exact-$k$ budget arm and the full 101{,}763-encounter readmission cohort,
and re-derived every number in the manuscript. \textbf{The screening criterion did not survive}: the
revised paper states that no dataset property measured here certifies a budget in advance, and that
the audit itself is the certificate. I also withdrew, before submission, an exploratory claim that
signal concentration orders the harm across high-dimensional datasets --- it failed when tested
across all rankers and statistics --- and I report a pre-registered causal prediction of my own that
the data refuted. The response letter details all three.

The revision also strengthens the central claim rather than only defending it. A
probability-independent confidence, computed from feature-space distance, shows that roughly
two-thirds of the measured selective-reliability harm reflects the probability scale while a third
persists; and a five-recalibrator flexibility ladder establishes that the failure of post-hoc
recalibration to recover selective reliability is not a limitation of the recalibrator.

\textbf{Fit and reproducibility.} The work is methodologically rigorous and fully reproducible: every
reported number regenerates from committed code, a hash-pinned calibration configuration and
recorded random seeds, and the repository is made public on the day of this resubmission, as
Reviewer~1 requested. The study uses only publicly available, de-identified datasets and adheres to
the TRIPOD+AI reporting guideline. Its question is of broad, practical relevance to machine
learning, clinical informatics and the responsible deployment of predictive models.

This manuscript is original, has not been published previously and is not under consideration
elsewhere. \textbf{It has a single author}, who has approved this submission, declares no competing
interests, and has had no prior discussions with a \textit{Scientific Reports} Editorial Board
Member about this work. Supplementary Information accompanies the submission.

Thank you for considering the revised manuscript.

\closing{Sincerely,}
\end{letter}
\end{document}
